\section{Team and competencies}
\label{sec:team}

\begin{evfutbox}
  No team exists yet. There is no host institution, no named principal investigator (PI) record
  for this program, no statistician, no CTA collaborator, no gate reader and no letter of support.
  The configurations below are planning constructs, sized from the PoC's own operating experience
  and from the work-package plan in \cref{sec:workpackages}. Every full-time-equivalent (FTE) figure is an indicative
  range, not a hire request.
\end{evfutbox}

\subsection{Team status today}

The PoC was designed, built, run and read by a single operator with AI assistance
(\cref{sec:threats-ai}). The same person was its only coder and is the prospective applicant; the
single-operator threat is rated in \cref{sec:threats}. This report names no other person and
claims no commitment from anyone. The following must be named, with written commitments, before
the step shown:

\begin{itemize}[noitemsep]
  \item \textbf{Before MS1 is read (Phase~0):} the non-owner blind coder, and the external gate
    reader.
  \item \textbf{Before the Phase~1 proposal is submitted:} the host institution and its
    person-month rate; the PI and the PI's record; a statistician as co-applicant; a CTA
    collaborator who can help locate itemized golds; a research engineer; a data-protection and
    license advisor.
  \item \textbf{Before Phase~2 is released:} a cognitive psychologist or CTA lead for V5; expert
    pools for the two held-out domains, with letters from professional bodies; a learning
    scientist.
  \item \textbf{Before V7 is scheduled (after MS5):} a partner's letter of intent.
\end{itemize}

\subsection{Why the program is interdisciplinary}

Each hard problem in the program belongs to a different research tradition, and no single skill
covers all of them.

\begin{itemize}[noitemsep]
  \item \textbf{The absence check} is a natural-language-processing and measurement problem
    (\cref{sec:n3,sec:resources}).
  \item \textbf{Eliciting a fair question from a predicted gap} is a cognitive-task-analysis
    problem, not a prompting technique (\cref{sec:foundations}).
  \item \textbf{Matching a predicted gap to what a novice finds hard} is a learning-science
    problem, because expert and instructor judgment of learner difficulty is unreliable
    (\cref{sec:education-blindspots}).
  \item \textbf{Deciding whether an effect is real at a pre-set size} is a statistics problem
    (\cref{sec:validation}).
  \item \textbf{Deploying on an organization's documents} is a security, privacy and legal problem
    (\cref{sec:organizations}).
\end{itemize}

\evint{} A team of machine-learning engineers alone cannot cover this list, and a team of domain
researchers alone cannot build the pipeline. The configurations below follow this map.

\subsection{Configurations by phase}

The team grows only when a gate releases the next phase.

\paragraph{Phase~0 (owner-run, no hires).} The owner prepares the units and the codebook. One
non-owner blind coder codes on an hourly contract, and the owner codes only under reported origin
blinding. The external gate reader reads MS1. Nothing else is needed.

\paragraph{Phase~1 lean team.} This is the minimum team that removes the PoC's largest structural
risk, one person who built, ran, coded and read everything, and that can reach MS3
(\cref{tab:team-lean}). It funds the area score and freeze (WP2), the judge's text-type check
(WP3), and the pooled retrospective study (WP4). The Git adapter and V3 are no longer in this
phase.

\begin{table}[htbp]
  \centering
  \footnotesize
  \caption{Phase~1 lean team (12--15 months, about 2.5--3.0~FTE, 35--44 person-months). It
    reaches the screening gate MS3.}
  \label{tab:team-lean}
  \begin{tabularx}{\linewidth}{@{}p{0.19\linewidth}Xp{0.13\linewidth}p{0.07\linewidth}p{0.12\linewidth}@{}}
    \toprule
    Role & Why needed & When (WP) & FTE & Permanent / consultant \\
    \midrule
    PI / research lead &
      Owns hypotheses and thresholds; prepares each gate report, but does not read the gate &
      WP1, all & 0.8--1.0 & Permanent \\
    Research engineer (LLM/retrieval) &
      Area score, admissibility as code, dated-corpus mode, the freeze, V0 &
      WP2, WP4 & 1.0 & Permanent \\
    Research assistant or second engineer &
      Gold availability check, area partitions, memorization probes &
      WP4 & 0.5 & Fixed-term \\
    Statistician / methodologist &
      Threshold table, pooled random-effects $\Delta$AUROC, incremental analysis; script
      committed before any gold is seen &
      WP1, WP3, WP4 & 0.3 & Part-time \\
    CTA practitioner &
      Codebook upkeep; rules for mapping gold items to areas &
      WP3, WP4 & 0.1--0.2 & Consultant \\
    Coders (2) &
      Judge text-type units; blind gold-to-area mapping &
      WP3, WP4 & 160--380 hours & Hourly contract \\
    Data-protection / license advisor &
      Dataset and gold license checks &
      WP1, ad hoc & Ad hoc & Consultant \\
    External gate reader &
      Reads MS3; not a team member (\cref{sec:gate-reader}) &
      MS3 & Per gate & Independent \\
    \bottomrule
  \end{tabularx}
\end{table}

\paragraph{Phase~2 team.} \evint{} This is the smallest team that can run Phase~2 without leaving
any stage's human-science, statistics or engineering skill missing when it starts
(\cref{tab:team-grant}). It adds a cognitive psychologist, a platform engineer, a learning
scientist, human--computer interaction (HCI) support, and the legal and project-management capacity V7 needs. Without the
organizational track, the platform engineer and legal counsel lines can be dropped. The role
table is in \cref{app:program}, since it matters only once MS3 reads continue.


\paragraph{Product / platform team.} \evfut{} This configuration starts only after MS5 and MS7
both read continue (\cref{sec:workpackages}). A research grant does not fund it; it is sized only
for comparison with the research teams above, at close to the \enquote{team of 8--15} used in
\cref{sec:resources}. The role table is \cref{tab:team-product} in \cref{app:program}.

\subsection{The external gate reader}
\label{sec:gate-reader}

\evint{} A funder buying pre-registered kill criteria needs someone other than the applicant to
apply them. The program therefore names an external gate reader: an independent statistician, or
a small advisory board that includes one.

\begin{itemize}[noitemsep]
  \item \textbf{Independence.} The reader is not a co-applicant, is not paid from a tranche that
    the reading releases, and has no stake in the program's continuation. A per-gate fee is agreed
    in advance.
  \item \textbf{Authority.} The reader reads MS1, MS3 and MS5 against the pre-registration,
    which is hash-committed in advance. The reading is binding: a tranche is released only if the reader reads
    continue.
  \item \textbf{Owner exceptions.} Any exception the owner wants at any gate, including MS2, MS2b,
    MS4, MS6, MS7, MS8 and the H-OSS gate, needs the reader's written sign-off before it takes effect. An exception
    without sign-off voids the gate reading, and no money is released on it.
\end{itemize}

No such reader existed during the PoC. The step from N2 to N3 was an owner exception, recorded
but not checked by anyone independent (\cref{sec:evolution}). The gate reader exists so that this
cannot recur.

\subsection{The independence rule (binding for staffing)}

One rule constrains every configuration above and cannot be relaxed by adding headcount. A domain
expert who serves as a study participant cannot also be an advisor who has seen the
reconstruction or its candidates for that domain. The same rule already governs the frozen
Validation Track~A. There, anyone who has seen reconstructed conservation content is ineligible
for any Track~A human role: interviewer, corroborating coder, K1 \enquote{absent} coder, or
importance rater (\emph{REORIENTATION.md} \S18.1). The NOW program extends the same logic to every
new domain:

\begin{itemize}[noitemsep]
  \item \textbf{Advisor versus participant.} An advisor who saw a domain's gap map or candidates
    may not later rate, be interviewed, code, or corroborate in that domain (eligibility register,
    WP1).
  \item \textbf{The Track~A roster stays separate.} Mechanics-domain staff are excluded from any
    Track~A role; E-SEAL predictions are script-generated, never shown to a future Track~A holder.
  \item \textbf{The pipeline's owner cannot be the blind coder.} Every gating agreement score includes a
    non-owner coder; the owner may code only under reported origin blinding, as in Phase~0.
  \item \textbf{Interviewers and coders never see scores or strata.} A guess-the-stratum check on a
    fifth of V5's areas measures whether this blinding held.
\end{itemize}

\evint{} These rules cost recruiting effort. Each domain needs at least one advisor and a disjoint
pool of participants, so \cref{sec:resources} counts recruitment overhead as its own line.
